\chapter{Limitations and Future Work}
\label{chap:limitations}

This project is a deliberately scoped teaching implementation, not a production-grade server or client. This chapter records, honestly and specifically, what has been intentionally left out, so that a reader extending this code understands exactly what remains to be done and why each omission was made.

\section{Concurrency}

Both servers handle exactly one client connection at a time: the \inlinecode{accept()} loop in \inlinecode{main()} (Chapters~\ref{chap:gopher-server} and~\ref{chap:http-server}) calls \inlinecode{handle\_client()} synchronously, and does not call \inlinecode{accept()} again until that call returns. A second client attempting to connect while the first is being served will simply wait in the kernel's connection backlog queue (sized by the \inlinecode{backlog} argument passed to \inlinecode{listen()}) until the first client's request is fully handled.

For the small, fast, local file-serving workloads this project targets, this is rarely noticeable in practice. A production server needs some form of concurrency to serve multiple clients simultaneously. The two classical approaches on POSIX systems are:

\begin{itemize}
    \item \textbf{Process-per-connection}, using \inlinecode{fork()} immediately after \inlinecode{accept()} returns, letting the child process handle that one client while the parent immediately loops back to \inlinecode{accept()} the next.
    \item \textbf{Thread-per-connection}, using \inlinecode{pthread\_create()} similarly, which avoids the memory overhead of a full process fork at the cost of requiring care around any shared mutable state.
\end{itemize}

Either would be a natural first extension to this codebase, and would require no changes to \inlinecode{socket\_utils.c} or the protocol-parsing logic --- only to the accept loop in each server's \inlinecode{main()}.

\section{HTTP Persistent Connections (Keep-Alive)}

Every response includes \texttt{Connection: close} (Chapter~\ref{chap:http-server}), and the client always closes its connection after one exchange (Chapter~\ref{chap:http-client}). Real HTTP/1.1 servers default to \emph{keep-alive}: reusing a single TCP connection for multiple sequential requests, which avoids the overhead of a new TCP handshake per request. Supporting this would require the server's \inlinecode{handle\_client} to loop, reading and responding to multiple requests per connection until the client explicitly closes it or a timeout elapses, and would require the client to read responses using the length-then-body strategy described in Chapter~\ref{chap:http-protocol} rather than the simpler ``read until close'' approach it currently reuses from the Gopher client (as noted explicitly in Chapter~\ref{chap:http-client}).

\section{Chunked Transfer Encoding}

This implementation always sends a fully-known \inlinecode{Content-Length}, computed by reading an entire file into memory before responding (Chapter~\ref{chap:http-server}). Real HTTP servers often need to stream responses whose total length is not known in advance (e.g.\ dynamically generated content), for which HTTP defines \textbf{chunked transfer encoding} as an alternative framing mechanism. Implementing it would extend the framing discussion in Chapter~\ref{chap:http-protocol} with a second, more complex, self-delimiting format.

\section{Robustness of Request Parsing}

Both servers read an incoming request with a single \inlinecode{recv()} call (Chapters~\ref{chap:gopher-server} and~\ref{chap:http-server}), rather than looping until a complete request is confirmed present, as the general principle in Section~\ref{sec:partial-io} would suggest for full robustness. This works reliably for the small requests generated by this project's own clients on the loopback interface, but is not guaranteed correct in general --- a sufficiently large or awkwardly-fragmented request could in principle arrive split across multiple TCP segments and thus multiple \inlinecode{recv()} calls. A fully robust server would apply the same ``read headers, find the boundary, then read exactly \inlinecode{Content-Length} more'' logic on the \emph{server} side that this project's HTTP \emph{client} already implements (Chapter~\ref{chap:http-client}).

\section{Gopher Path-Traversal Protection}

As noted explicitly in Chapter~\ref{chap:gopher-server}, \inlinecode{send\_file} in \texttt{gopher\_server.c} performs no validation of the requested selector before using it to construct a filesystem path, unlike the HTTP server's deliberate (if simple) \texttt{".."} check. Bringing the same defense to the Gopher server --- or, better, replacing both implementations' substring check with proper path canonicalization and a containment check against \inlinecode{root\_dir} --- would be a natural security-hardening exercise.

\section{Input Validation}

Command-line port arguments are parsed with \inlinecode{atoi()} in both clients (Chapters~\ref{chap:gopher-client} and~\ref{chap:http-client}), which silently returns \texttt{0} for non-numeric input rather than reporting an error. Using \inlinecode{strtol()} with proper error checking (verifying the entire string was consumed and the result falls within the valid port range 1--65535) would make failures far more diagnosable.

\section{HTTP Method Support}

Only \texttt{GET} is implemented; every other method receives \texttt{405 Method Not Allowed} (Chapter~\ref{chap:http-server}). Adding \texttt{POST} support would introduce the need to parse a request body, which in turn requires applying the same header/body framing logic from Chapter~\ref{chap:http-protocol} on the \emph{server} side --- currently the server never needs to look past the blank line terminating the headers, since GET requests carry no body.

\section{Summary}

None of the above omissions are accidental; each represents a conscious boundary drawn to keep this project's scope focused on the core networking and framing concepts documented in Chapters~\ref{chap:networking-fundamentals} through~\ref{chap:http-client}. Readers looking to extend this project are encouraged to treat this chapter as a roadmap, tackling one item at a time against the existing, tested foundation described in Chapter~\ref{chap:testing}.
